\label{chap:Introduction}

Catalysis research generates complex and heterogeneous experimental data.
A single dataset may contain information about catalyst materials, synthesis and treatment steps, reactor configurations, operating conditions, analytical measurements, processing parameters, and performance results.
These elements become scientifically meaningful only when their experimental context is preserved.
For example, a conversion value, an \ac{NMR} spectrum, or a kinetic measurement cannot be interpreted reliably without information about the investigated material, the sample preparation, the applied method, and the conditions under which the data were produced \cite{mendes2021opendata,mendes2025data}.

This need for contextual information is directly connected to the \ac{FAIR} principles.
Catalysis data should not only be stored and made accessible, but also be described in a way that allows researchers and machines to find, interpret, compare, and reuse them \cite{wilkinson2016fair,mendes2025data}.
However, much of the relevant metadata is still embedded in heterogeneous file structures, instrument exports, tables, laboratory notes, captions, or human-readable descriptions.
As a result, experimental information may technically be available, while still being difficult to transform into machine-actionable metadata.

Semantic metadata standards and controlled vocabularies offer a way to address this problem.
Application profiles such as \ac{DCAT-AP+} and \ac{ChemDCAT-AP} define structured ways to describe datasets, activities, entities, attributes, and provenance relations \cite{stroemert2026chemdcat,dcatapplusDocumentation,chemdcatapDocumentation}.
Vocabulary resources such as \ac{Voc4Cat}, \ac{QUDT}, and \ac{nmrCV} provide controlled identifiers for catalysis concepts, quantities, units, and \ac{NMR}-specific metadata \cite{david_linke_2023_8313340,qudt2022schema,Schober2018nmrML}.
Together, these resources can support interoperable and reusable data descriptions.
Nevertheless, they do not automatically extract the required information from real experimental dataset packages.

This creates an operational gap between available experimental data and semantically structured metadata records.
Manual curation can close this gap, but it is time-consuming, difficult to scale, and dependent on domain expertise.
Large Language Models (LLMs) offer a promising approach for supporting this process because they can interpret heterogeneous textual sources and generate structured outputs \cite{dagdelen2024structured,schilling2025text,xu2024generativeie}.
However, automatic extraction alone is not sufficient.
The resulting metadata must remain traceable to source evidence, validated against a target structure, grounded in controlled vocabularies, and inspectable by human curators.

This thesis is motivated by the need for a workflow that connects these layers.
It investigates how an LLM-supported process can extract metadata signals from heterogeneous catalytic experiment datasets and organize them into an evidence-backed semantic metadata draft.
The goal is not to replace expert curation, but to reduce the manual effort required to produce structured, \ac{FAIR}-oriented metadata and to make the extraction process more transparent, reproducible, and semantically grounded.
In this work, \ac{NMR} datasets serve as a test case because they combine raw instrument data, processing parameters, standardized exchange formats, and chemically meaningful spectral information within a compact but realistic experimental package \cite{Schober2018nmrML}.
